\scanpage{69}
\mspage{\handwritten{- 4 -}}

p.37, li.6 : ``composé'' au lieu de ``composée''.

\begin{center}
Commentaires Appendice
\end{center}

p.3 : détailler preuve de 7.4 ? ($Y$ plat sur $S$, ... il en résulte aisément ..., on vérifie aisément ...)

p.1 : dans l'énoncé 7.1, de quel $K^{\bullet}$ s'agit-il ?

p.5, li.5 : réf. pour \handwritten{``}$X_t$ et $X_s$ réguliers $\handwritten{\Rightarrow}$\handwritten{\unsure{``|''}}$X$ régulier\handwritten{''}.

p.6, li.6 : peut-on voir l'argument ?

p.7, li.8 du bas : ``Ce résultat généralise 7.5'' (une fois qu'on a remarqué que, pour $X = S$, les flèches du diagramme de 7.3 sont les flèches identiques de $\underline{Z}$.)

p.7, li.5 du bas : ``récrire'' et non ``réécrire''.

p.8, li.12 du bas : ``horizontales'' et non ``verticales''.

p.8, li.6 du bas : ``du prisme'' et non ``du carré''.

p.9 : nommer les flèches verticales d\deleted{a}e (7.10.2 bis).

p.9, li.6 : ajouter ``(cf. 3.3)'' après \unsure{```|``}implique''.

% A small detached, raised quadrilateral follows the question mark.
% Its status as authored punctuation rather than a surface speck is uncertain.
% Inline TikZ preserves its position without assigning it a punctuation value.
p.9 : références pour $\mathrm{Gr}^{\bullet}_{\mathrm{top}}$ et l'assertion de compatibilité aux filtrations ?\unsure{\begin{tikzpicture}[x=1mm,y=1mm,baseline=0pt]
\fill (0.65,0.80) -- (0.91,0.80) -- (1.02,1.08) -- (0.76,1.17) -- cycle;
\end{tikzpicture}}

% Two detached dark marks lie below the r / raised point of the second Gr.
% Their punctuation/accidental status is unresolved; their shapes are retained.
Donner aussi le diagramme analogue à (7.10.1 bis) avec \makebox[0pt][l]{\hspace{0.90em}\raisebox{-1.45ex}[0pt][0pt]{\unsure{\begin{tikzpicture}[x=1mm,y=1mm,baseline=0pt]
\fill (0,0) -- (0.22,0.09) -- (0.04,0.24) -- (-0.04,0.17) -- cycle;
\fill (0.33,0.48) .. controls (0.46,0.52) and (0.58,0.73) .. (0.49,0.78) .. controls (0.36,0.78) and (0.22,0.57) .. (0.33,0.48) -- cycle;
\end{tikzpicture}}}}$\mathrm{Gr}^{\bullet}(\ )_{\underline{Q}}$ pour $g$ projectif, d'intersection complète, et de dimension relative virtuelle donnée.

p.10, li.11 : ``standard'' et non``standart''.

p.10, li.13 : ne faut-il pas ajouter ``connexe'' après ``schéma algébrique'' et ``courbe algébrique'' ? Y a-t-il une référence à l'intention du lecteur qui aimerait voir une démonstration d\deleted{d}e l'assertion en question?

% archive sheet number: 68
